\section{Movimiento en dos y tres dimensiones}

Hasta ahora hemos estudiado principalmente movimientos que pueden describirse
mediante una sola coordenada. Por ejemplo, un automóvil que avanza por una
carretera completamente recta puede describirse, durante cierto intervalo,
mediante una posición \(x(t)\).

Sin embargo, una gran cantidad de movimientos reales no ocurre sobre una sola
línea.

Un automóvil puede tomar una curva, un dron puede desplazarse simultáneamente
hacia adelante y hacia arriba, una embarcación puede cambiar su posición en dos
direcciones sobre la superficie del agua y una aeronave puede modificar las tres
coordenadas espaciales durante su movimiento.

En estas situaciones ya no basta una sola coordenada.

Necesitamos describir la posición, el desplazamiento, la velocidad y la
aceleración mediante \textbf{vectores}.

\begin{idea}[La extensión fundamental]
El paso del movimiento en una dimensión al movimiento en dos o tres dimensiones
no exige abandonar los conceptos de cinemática ya construidos.

Las ideas siguen siendo las mismas:

\[
\text{posición}
\longrightarrow
\text{desplazamiento}
\longrightarrow
\text{velocidad}
\longrightarrow
\text{aceleración}.
\]

La diferencia fundamental es que estas magnitudes deben tratarse ahora como
vectores y, por tanto, poseen componentes, magnitud, dirección y sentido.
\end{idea}


\section{¿Cuántas dimensiones necesita un movimiento?}

La cantidad de dimensiones necesarias depende de cuántas coordenadas espaciales
deben cambiar para describir el movimiento.

\subsection{Movimiento en una dimensión}

Si el objeto permanece sobre una línea recta, una sola coordenada puede ser
suficiente:

\[
x=x(t).
\]

Por ejemplo, el movimiento de un ascensor idealizado puede describirse utilizando
únicamente una coordenada vertical.

\subsection{Movimiento en dos dimensiones}

Si la partícula se mueve dentro de un plano, necesitamos dos coordenadas:

\[
x=x(t),
\qquad
y=y(t).
\]

El movimiento ocurre en el plano \(xy\).

\subsection{Movimiento en tres dimensiones}

Para describir un movimiento espacial general necesitamos tres coordenadas:

\[
x=x(t),
\qquad
y=y(t),
\qquad
z=z(t).
\]

\begin{interpretacion}
La dimensionalidad del movimiento no depende de que el objeto real sea
tridimensional.

Un automóvil es un cuerpo tridimensional, pero si se mueve a lo largo de una
carretera recta y sólo queremos estudiar su traslación, su movimiento puede
modelarse como unidimensional.

La dimensionalidad se refiere al número de coordenadas necesarias para describir
el movimiento que estamos estudiando.
\end{interpretacion}


\section{Modelo de partícula}

Continuaremos utilizando el modelo de partícula desarrollado anteriormente.

\begin{modelo}[Partícula]
Representamos el objeto mediante un punto cuando sus dimensiones, forma y
orientación no son importantes para responder la pregunta física planteada.

El punto conserva la información necesaria sobre la posición y el movimiento
traslacional del objeto, pero ignora detalles como su tamaño y su rotación.
\end{modelo}

Este modelo permite estudiar mediante las mismas ecuaciones el movimiento de
objetos físicamente muy diferentes, siempre que sus dimensiones no sean
relevantes para el problema considerado.


\section{Sistema de referencia en dos dimensiones}

Antes de describir un movimiento debemos establecer un sistema de referencia.

En un plano cartesiano utilizamos dos ejes mutuamente perpendiculares:

\[
x
\qquad\text{y}\qquad
y.
\]

Definimos también los vectores unitarios:

\[
\ihat
\qquad\text{y}\qquad
\jhat,
\]

que apuntan, respectivamente, hacia los sentidos positivos de los ejes \(x\) y
\(y\).

\begin{center}
\begin{tikzpicture}[scale=1.0]

    \draw[-{Latex},thick] (-0.5,0) -- (5.3,0)
        node[right] {\(x\)};

    \draw[-{Latex},thick] (0,-0.5) -- (0,4.2)
        node[above] {\(y\)};

    \fill (3.5,2.6) circle (2.5pt);
    \node[above right] at (3.5,2.6) {\(P(x,y)\)};

    \draw[dashed] (3.5,0) -- (3.5,2.6);
    \draw[dashed] (0,2.6) -- (3.5,2.6);

    \draw[-{Latex},very thick,AzulPrincipal]
        (0,0) -- (3.5,2.6)
        node[midway,above left] {\(\vect{r}\)};

    \node[below] at (3.5,0) {\(x\)};
    \node[left] at (0,2.6) {\(y\)};

    \draw[-{Latex},thick,VerdePrincipal]
        (0.4,0.45) -- (1.25,0.45)
        node[midway,below] {\(\ihat\)};

    \draw[-{Latex},thick,VerdePrincipal]
        (0.4,0.45) -- (0.4,1.30)
        node[midway,left] {\(\jhat\)};

\end{tikzpicture}
\end{center}

\begin{advertencia}
Los signos de las componentes \(x\) y \(y\) no poseen significado absoluto.

Su interpretación depende de:

\begin{itemize}
    \item dónde se encuentra el origen;
    \item cómo están orientados los ejes;
    \item qué sentidos se eligieron como positivos.
\end{itemize}

Estas decisiones deben establecerse antes de analizar el movimiento.
\end{advertencia}


\section{Sistema de referencia en tres dimensiones}

Para un movimiento espacial introducimos un tercer eje:

\[
z,
\]

perpendicular a los otros dos, y el vector unitario correspondiente:

\[
\khat.
\]

Utilizaremos un sistema cartesiano de mano derecha formado por:

\[
\ihat,
\qquad
\jhat,
\qquad
\khat.
\]

Una posición espacial se especificará entonces mediante las tres coordenadas:

\[
(x,y,z).
\]

\begin{recordatorio}[Componentes rectangulares]
En todo este capítulo utilizaremos prioritariamente componentes cartesianas.

Para cualquier vector tridimensional:

\[
\vect{A}
=
A_x\ihat+A_y\jhat+A_z\khat.
\]

Su magnitud es:

\[
A
=
\norm{\vect{A}}
=
\sqrt{A_x^2+A_y^2+A_z^2}.
\]

En dos dimensiones basta tomar:

\[
A_z=0.
\]
\end{recordatorio}


\section{Vector de posición}

En movimiento rectilíneo utilizábamos una coordenada \(x\) para indicar dónde se
encontraba la partícula.

En dos o tres dimensiones necesitamos un vector.

\begin{definicion}[Vector de posición]
El vector de posición \(\vect{r}\) es el vector que se extiende desde el origen
del sistema de referencia hasta la posición ocupada por la partícula.
\end{definicion}

En tres dimensiones:

\[
\boxed{
\vect{r}
=
x\ihat+y\jhat+z\khat
}.
\]

En dos dimensiones:

\[
\boxed{
\vect{r}
=
x\ihat+y\jhat
}.
\]

Las coordenadas:

\[
x,\qquad y,\qquad z
\]

son las componentes rectangulares del vector de posición.

En el SI se expresan en metros.

\subsection{Magnitud del vector de posición}

La magnitud del vector de posición es:

\[
r
=
\norm{\vect{r}}
=
\sqrt{x^2+y^2+z^2}.
\]

En un movimiento bidimensional:

\[
r=\sqrt{x^2+y^2}.
\]

Esta magnitud representa la distancia geométrica entre el origen y la posición
actual de la partícula.

\begin{errorfrecuente}[Magnitud de la posición y distancia recorrida]
La cantidad:

\[
\norm{\vect{r}}
\]

no representa, en general, la distancia que ha recorrido la partícula.

Solamente indica qué tan lejos se encuentra la posición actual respecto del
origen.

Una partícula podría haber seguido una trayectoria muy larga y terminar cerca
del origen.
\end{errorfrecuente}


\section{La posición depende del tiempo}

Si una partícula se mueve, sus coordenadas pueden cambiar con el tiempo:

\[
x=x(t),
\qquad
y=y(t),
\qquad
z=z(t).
\]

Por tanto, podemos escribir:

\[
\boxed{
\vect{r}(t)
=
x(t)\ihat+y(t)\jhat+z(t)\khat
}.
\]

Esta función vectorial contiene la descripción cinemática de la posición de la
partícula en cada instante.

\begin{interpretacion}
Conocer \(\vect{r}(t)\) significa conocer simultáneamente:

\[
x(t),\qquad y(t),\qquad z(t).
\]

No son tres tiempos diferentes ni tres movimientos desconectados.

Son tres componentes espaciales del mismo movimiento y todas se evalúan en el
mismo instante \(t\).
\end{interpretacion}


\section{Trayectoria}

A medida que cambia el vector de posición, la partícula va ocupando diferentes
puntos del espacio.

El conjunto de esos puntos forma su trayectoria.

\begin{definicion}[Trayectoria]
La trayectoria es la curva geométrica formada por las posiciones sucesivas que
ocupa una partícula durante su movimiento.
\end{definicion}

Una trayectoria puede ser:

\begin{itemize}
    \item rectilínea;
    \item curvilínea;
    \item plana;
    \item espacial.
\end{itemize}

\begin{center}
\begin{tikzpicture}[scale=0.95]

    \draw[-{Latex},thick] (-0.5,0) -- (6.6,0)
        node[right] {\(x\)};

    \draw[-{Latex},thick] (0,-0.5) -- (0,4.5)
        node[above] {\(y\)};

    \draw[
        very thick,
        AzulSecundario,
        smooth
    ]
    plot coordinates {
        (0.7,0.8)
        (1.6,1.7)
        (2.6,2.2)
        (3.8,2.0)
        (4.8,2.8)
        (5.8,3.7)
    };

    \fill (0.7,0.8) circle (2.5pt);
    \node[below right] at (0.7,0.8) {\(P_i\)};

    \fill (5.8,3.7) circle (2.5pt);
    \node[above left] at (5.8,3.7) {\(P_f\)};

    \node[above] at (3.1,2.45) {trayectoria};

\end{tikzpicture}
\end{center}


\section{Ecuaciones paramétricas de la trayectoria}

Cuando conocemos:

\[
x=x(t),
\qquad
y=y(t),
\]

estas dos ecuaciones describen la trayectoria de manera paramétrica.

El tiempo \(t\) actúa como parámetro común.

En tres dimensiones tendríamos:

\[
x=x(t),
\qquad
y=y(t),
\qquad
z=z(t).
\]

\subsection{Eliminar el tiempo}

En algunos casos podemos eliminar \(t\) entre las ecuaciones y obtener una
relación directa entre las coordenadas.

Por ejemplo, si:

\[
x=2t
\]

y:

\[
y=t^2,
\]

de la primera relación:

\[
t=\frac{x}{2}.
\]

Sustituyendo en la segunda:

\[
y=\left(\frac{x}{2}\right)^2.
\]

Por tanto:

\[
\boxed{
y=\frac{x^2}{4}
}.
\]

Esta ecuación describe la forma geométrica de la trayectoria en el plano \(xy\).

\begin{interpretacion}
Las ecuaciones:

\[
x(t),\qquad y(t)
\]

indican \textbf{cómo cambia la posición con el tiempo}.

La ecuación:

\[
y=y(x)
\]

describe \textbf{la forma geométrica de la trayectoria} después de haber
eliminado el tiempo.

Son representaciones diferentes del mismo movimiento.
\end{interpretacion}


\section{Desplazamiento en dos y tres dimensiones}

Supongamos que en el instante \(t_i\) una partícula tiene posición:

\[
\vect{r}_i,
\]

y en un instante posterior \(t_f\):

\[
\vect{r}_f.
\]

El desplazamiento continúa siendo el cambio de posición.

\begin{definicion}[Vector desplazamiento]
El desplazamiento de una partícula entre dos posiciones es:

\[
\boxed{
\Delta\vect{r}
=
\vect{r}_f-\vect{r}_i
}.
\]
\end{definicion}

Si:

\[
\vect{r}_i
=
x_i\ihat+y_i\jhat+z_i\khat
\]

y:

\[
\vect{r}_f
=
x_f\ihat+y_f\jhat+z_f\khat,
\]

entonces:

\[
\Delta\vect{r}
=
(x_f-x_i)\ihat
+
(y_f-y_i)\jhat
+
(z_f-z_i)\khat.
\]

Definiendo:

\[
\Delta x=x_f-x_i,
\]

\[
\Delta y=y_f-y_i,
\]

\[
\Delta z=z_f-z_i,
\]

obtenemos:

\[
\boxed{
\Delta\vect{r}
=
\Delta x\,\ihat
+
\Delta y\,\jhat
+
\Delta z\,\khat
}.
\]


\subsection{Interpretación geométrica}

El vector desplazamiento apunta directamente desde la posición inicial hasta la
posición final.

\begin{center}
\begin{tikzpicture}[scale=0.95]

    \draw[-{Latex},thick] (-0.5,0) -- (6.5,0)
        node[right] {\(x\)};

    \draw[-{Latex},thick] (0,-0.5) -- (0,4.6)
        node[above] {\(y\)};

    \coordinate (A) at (1.2,1.0);
    \coordinate (B) at (5.3,3.4);

    \draw[
        very thick,
        AzulSecundario,
        smooth
    ]
    plot coordinates {
        (1.2,1.0)
        (2.0,2.0)
        (3.1,2.6)
        (4.1,2.2)
        (5.3,3.4)
    };

    \draw[-{Latex},thick,AzulPrincipal]
        (0,0) -- (A)
        node[midway,below right] {\(\vect{r}_i\)};

    \draw[-{Latex},thick,AzulPrincipal]
        (0,0) -- (B)
        node[midway,above left] {\(\vect{r}_f\)};

    \draw[-{Latex},very thick,VerdePrincipal]
        (A) -- (B)
        node[midway,below right] {\(\Delta\vect{r}\)};

    \fill (A) circle (2.5pt);
    \fill (B) circle (2.5pt);

    \node[left] at (A) {\(P_i\)};
    \node[right] at (B) {\(P_f\)};

\end{tikzpicture}
\end{center}

\begin{interpretacion}
El desplazamiento no describe todo el camino recorrido.

Sólo compara las posiciones inicial y final.

Por ello, dos trayectorias completamente diferentes pueden producir exactamente
el mismo vector desplazamiento si comienzan y terminan en los mismos puntos.
\end{interpretacion}


\section{Distancia recorrida y magnitud del desplazamiento}

La distancia recorrida mide la longitud total de la trayectoria.

La magnitud del desplazamiento es:

\[
\norm{\Delta\vect{r}}
=
\sqrt{
(\Delta x)^2
+
(\Delta y)^2
+
(\Delta z)^2
}.
\]

En general:

\[
\boxed{
d\geq\norm{\Delta\vect{r}}
}.
\]

La igualdad se produce solamente cuando la trayectoria recorrida entre las dos
posiciones coincide con el segmento recto que las une y no existe retroceso.

Para una descripción continua, la distancia recorrida puede escribirse como:

\[
d
=
\int_{t_i}^{t_f}
\norm{\vect{v}(t)}
\,\dd t.
\]

Más adelante veremos que:

\[
\norm{\vect{v}(t)}
\]

es precisamente la rapidez instantánea.

\begin{errorfrecuente}
No debe calcularse la distancia total simplemente mediante:

\[
\sqrt{(\Delta x)^2+(\Delta y)^2}
\]

cuando la trayectoria es curva.

Esa expresión proporciona la magnitud del \textbf{desplazamiento}, no la longitud
de una trayectoria arbitraria.
\end{errorfrecuente}


\section{Ejemplo: posición y desplazamiento en un plano}

\begin{ejemplo}[Movimiento entre dos puntos]
Una partícula se mueve en un plano desde:

\[
P_i=(-2,1)\,\si{\meter}
\]

hasta:

\[
P_f=(4,5)\,\si{\meter}.
\]

Determine el vector desplazamiento y su magnitud.
\end{ejemplo}

Los vectores de posición son:

\[
\vect{r}_i
=
(-2\ihat+\jhat)\,\si{\meter},
\]

\[
\vect{r}_f
=
(4\ihat+5\jhat)\,\si{\meter}.
\]

El desplazamiento es:

\[
\Delta\vect{r}
=
\vect{r}_f-\vect{r}_i.
\]

Trabajando por componentes:

\[
\Delta x
=
4-(-2)
=
\SI{6}{\meter},
\]

\[
\Delta y
=
5-1
=
\SI{4}{\meter}.
\]

Por tanto:

\[
\boxed{
\Delta\vect{r}
=
(6\ihat+4\jhat)\,\si{\meter}
}.
\]

Su magnitud es:

\[
\norm{\Delta\vect{r}}
=
\sqrt{6^2+4^2}\,\si{\meter}.
\]

Así:

\[
\boxed{
\norm{\Delta\vect{r}}
=
\sqrt{52}\,\si{\meter}
\approx
\SI{7.21}{\meter}
}.
\]

\begin{interpretacion}
El resultado indica que la posición final está:

\begin{itemize}
    \item \(\SI{6}{\meter}\) más hacia \(+x\);
    \item \(\SI{4}{\meter}\) más hacia \(+y\);
\end{itemize}

que la posición inicial.

La magnitud \(\SI{7.21}{\meter}\) es la separación recta entre las posiciones
inicial y final. No podemos determinar la distancia realmente recorrida sin
conocer la trayectoria.
\end{interpretacion}


\section{Velocidad media}

En una dimensión definimos la velocidad media como el desplazamiento dividido
entre el intervalo de tiempo.

La misma idea se conserva en varias dimensiones.

\begin{definicion}[Velocidad media]
La velocidad media de una partícula durante un intervalo \(\Delta t\) es:

\[
\boxed{
\vect{\bar{v}}
=
\frac{\Delta\vect{r}}{\Delta t}
}.
\]
\end{definicion}

Como:

\[
\Delta t=t_f-t_i>0,
\]

la velocidad media tiene la misma dirección y sentido que el desplazamiento.

En componentes:

\[
\boxed{
\bar{v}_x
=
\frac{\Delta x}{\Delta t}
},
\]

\[
\boxed{
\bar{v}_y
=
\frac{\Delta y}{\Delta t}
},
\]

\[
\boxed{
\bar{v}_z
=
\frac{\Delta z}{\Delta t}
}.
\]

Por tanto:

\[
\vect{\bar{v}}
=
\bar{v}_x\ihat
+
\bar{v}_y\jhat
+
\bar{v}_z\khat.
\]

La unidad SI es:

\[
\si{\meter\per\second}.
\]

\begin{interpretacion}
La velocidad media contiene información sobre el cambio neto de posición por
unidad de tiempo.

No describe necesariamente cómo fue la velocidad durante cada instante del
recorrido.
\end{interpretacion}


\subsection{Ejemplo}

Si en el ejemplo anterior la partícula tarda:

\[
\Delta t=\SI{2.0}{\second},
\]

su velocidad media es:

\[
\vect{\bar{v}}
=
\frac{
(6\ihat+4\jhat)\,\si{\meter}
}{
\SI{2.0}{\second}
}.
\]

Por tanto:

\[
\boxed{
\vect{\bar{v}}
=
(3.0\ihat+2.0\jhat)\,
\si{\meter\per\second}
}.
\]

Su magnitud es:

\[
\norm{\vect{\bar{v}}}
=
\sqrt{3.0^2+2.0^2}\,
\si{\meter\per\second}.
\]

Así:

\[
\boxed{
\norm{\vect{\bar{v}}}
\approx
\SI{3.61}{\meter\per\second}
}.
\]


\section{Velocidad instantánea}

La velocidad media describe un intervalo completo.

Sin embargo, con frecuencia queremos saber cómo se mueve la partícula en un
instante determinado.

Para ello reducimos progresivamente el intervalo temporal.

\[
\vect{\bar{v}}
=
\frac{\Delta\vect{r}}{\Delta t}.
\]

Cuando:

\[
\Delta t\to0,
\]

el desplazamiento también se hace muy pequeño y obtenemos la velocidad
instantánea.

\begin{definicion}[Velocidad instantánea]
La velocidad instantánea es:

\[
\boxed{
\vect{v}
=
\lim_{\Delta t\to0}
\frac{\Delta\vect{r}}{\Delta t}
=
\deriv{\vect{r}}{t}
}.
\]
\end{definicion}

Esta ecuación significa físicamente que la velocidad es la tasa instantánea de
cambio del vector de posición con respecto al tiempo.


\section{La velocidad es tangente a la trayectoria}

Esta propiedad es una de las ideas geométricas más importantes del movimiento en
dos y tres dimensiones.

Durante un intervalo finito, el desplazamiento:

\[
\Delta\vect{r}
\]

une dos puntos de la trayectoria.

La velocidad media tiene la misma dirección que ese desplazamiento.

A medida que hacemos más pequeño el intervalo:

\[
\Delta t\to0,
\]

los dos puntos se aproximan y la dirección del desplazamiento se aproxima a la
dirección tangente a la curva.

Por ello:

\begin{principio}[Dirección de la velocidad instantánea]
En cada punto de una trayectoria, el vector velocidad instantánea es tangente a
la trayectoria y apunta en el sentido instantáneo del movimiento.
\end{principio}

\begin{center}
\begin{tikzpicture}[scale=1.0]

    \draw[-{Latex},thick] (-0.5,0) -- (6.6,0)
        node[right] {\(x\)};

    \draw[-{Latex},thick] (0,-0.5) -- (0,4.5)
        node[above] {\(y\)};

    \draw[
        very thick,
        AzulSecundario,
        smooth
    ]
    plot coordinates {
        (0.8,0.8)
        (1.7,1.4)
        (2.8,2.3)
        (3.7,2.9)
        (4.8,3.1)
        (5.8,2.7)
    };

    \coordinate (P) at (3.7,2.9);

    \fill (P) circle (2.5pt);
    \node[below right] at (P) {\(P\)};

    \draw[-{Latex},very thick,VerdePrincipal]
        (P) -- ++(1.7,0.45)
        node[right] {\(\vect{v}\)};

    \node[above left] at (2.6,2.25) {trayectoria};

\end{tikzpicture}
\end{center}

\begin{errorfrecuente}
La velocidad instantánea no apunta necesariamente desde el origen hacia la
partícula.

El vector:

\[
\vect{r}
\]

apunta desde el origen hasta la posición.

El vector:

\[
\vect{v}
\]

es tangente a la trayectoria.

Son vectores físicamente diferentes.
\end{errorfrecuente}


\section{Componentes de la velocidad}

Escribimos:

\[
\vect{r}
=
x\ihat+y\jhat+z\khat.
\]

Los vectores unitarios cartesianos son constantes cuando los ejes están fijos:

\[
\deriv{\ihat}{t}
=
\deriv{\jhat}{t}
=
\deriv{\khat}{t}
=
0.
\]

Por tanto:

\[
\vect{v}
=
\deriv{\vect{r}}{t}
=
\deriv{x}{t}\ihat
+
\deriv{y}{t}\jhat
+
\deriv{z}{t}\khat.
\]

Comparando con:

\[
\vect{v}
=
v_x\ihat+v_y\jhat+v_z\khat,
\]

obtenemos:

\[
\boxed{
v_x=\deriv{x}{t}
},
\]

\[
\boxed{
v_y=\deriv{y}{t}
},
\]

\[
\boxed{
v_z=\deriv{z}{t}
}.
\]

\begin{interpretacion}
Cada componente de la velocidad indica qué tan rápidamente cambia la coordenada
correspondiente.

Por ejemplo:

\[
v_x>0
\]

significa que, en ese instante, la coordenada \(x\) está aumentando.

En cambio:

\[
v_y<0
\]

significa que la coordenada \(y\) está disminuyendo.

Estos signos deben interpretarse respecto de los ejes previamente definidos.
\end{interpretacion}


\section{Rapidez}

La rapidez continúa siendo una magnitud escalar.

\begin{definicion}[Rapidez instantánea]
La rapidez es la magnitud del vector velocidad:

\[
\boxed{
v=\norm{\vect{v}}
}.
\]
\end{definicion}

En tres dimensiones:

\[
\boxed{
v
=
\sqrt{v_x^2+v_y^2+v_z^2}
}.
\]

En el plano:

\[
\boxed{
v
=
\sqrt{v_x^2+v_y^2}
}.
\]

La rapidez siempre es:

\[
v\geq0.
\]

\begin{advertencia}
No deben confundirse:

\[
\vect{v}
\]

y:

\[
v=\norm{\vect{v}}.
\]

La velocidad posee dirección y sentido.

La rapidez solamente representa su magnitud.
\end{advertencia}


\section{Dirección de la velocidad en un plano}

Si:

\[
\vect{v}
=
v_x\ihat+v_y\jhat,
\]

y medimos el ángulo \(\theta\) desde el eje \(+x\), entonces:

\[
v_x=v\cos\theta,
\]

\[
v_y=v\sin\theta.
\]

Cuando conocemos las componentes:

\[
\tan\theta=\frac{v_y}{v_x}.
\]

Así:

\[
\theta
=
\arctan\left(\frac{v_y}{v_x}\right),
\]

pero el cuadrante debe determinarse observando los signos de \(v_x\) y \(v_y\).

\begin{errorfrecuente}[El cuadrante]
La calculadora puede devolver un ángulo cuya tangente sea correcta pero que no
corresponda al cuadrante físico del vector.

Por ello, antes de dar una dirección final siempre deben examinarse los signos de:

\[
v_x
\qquad\text{y}\qquad
v_y.
\]
\end{errorfrecuente}


\section{Aceleración media}

La aceleración mide cómo cambia la velocidad.

Ahora debemos recordar que la velocidad es un vector.

Supongamos que:

\[
\vect{v}_i
\]

es la velocidad inicial y:

\[
\vect{v}_f
\]

la velocidad final.

El cambio de velocidad es:

\[
\Delta\vect{v}
=
\vect{v}_f-\vect{v}_i.
\]

\begin{definicion}[Aceleración media]
La aceleración media durante un intervalo \(\Delta t\) es:

\[
\boxed{
\vect{\bar{a}}
=
\frac{\Delta\vect{v}}{\Delta t}
}.
\]
\end{definicion}

En componentes:

\[
\bar{a}_x
=
\frac{\Delta v_x}{\Delta t},
\]

\[
\bar{a}_y
=
\frac{\Delta v_y}{\Delta t},
\]

\[
\bar{a}_z
=
\frac{\Delta v_z}{\Delta t}.
\]

La unidad SI es:

\[
\si{\meter\per\second\squared}.
\]


\section{Aceleración instantánea}

Al tomar el límite:

\[
\Delta t\to0,
\]

obtenemos la aceleración instantánea.

\begin{definicion}[Aceleración instantánea]
La aceleración instantánea es:

\[
\boxed{
\vect{a}
=
\lim_{\Delta t\to0}
\frac{\Delta\vect{v}}{\Delta t}
=
\deriv{\vect{v}}{t}
}.
\]
\end{definicion}

Como:

\[
\vect{v}
=
\deriv{\vect{r}}{t},
\]

también:

\[
\boxed{
\vect{a}
=
\derivdos{\vect{r}}{t}
}.
\]

Por componentes:

\[
\boxed{
a_x=\deriv{v_x}{t}
=
\derivdos{x}{t}
},
\]

\[
\boxed{
a_y=\deriv{v_y}{t}
=
\derivdos{y}{t}
},
\]

\[
\boxed{
a_z=\deriv{v_z}{t}
=
\derivdos{z}{t}
}.
\]


\section{¿Qué puede cambiar cuando existe aceleración?}

En movimiento rectilíneo era fácil concentrarse en el aumento o disminución de
la rapidez.

En varias dimensiones aparece una posibilidad adicional.

Como la velocidad es un vector, puede cambiar:

\begin{enumerate}
    \item su magnitud;
    \item su dirección;
    \item su magnitud y su dirección simultáneamente.
\end{enumerate}

Por tanto:

\begin{idea}[Acelerar no significa únicamente aumentar la rapidez]
Existe aceleración siempre que cambie el vector velocidad.

Eso puede ocurrir incluso si la rapidez permanece constante, siempre que cambie
la dirección del movimiento.
\end{idea}


\subsection{Cambio de magnitud}

Si el vector velocidad conserva su dirección pero aumenta o disminuye su
magnitud, existe aceleración.

Este caso incluye muchos movimientos rectilíneos.


\subsection{Cambio de dirección}

Supongamos que:

\[
\norm{\vect{v}_i}
=
\norm{\vect{v}_f},
\]

pero:

\[
\vect{v}_i\neq\vect{v}_f.
\]

La rapidez inicial y final pueden ser iguales, pero como cambió la dirección:

\[
\Delta\vect{v}\neq0.
\]

Por tanto:

\[
\vect{a}\neq0.
\]

\begin{errorfrecuente}[Rapidez constante no implica aceleración nula]
La afirmación:

\[
v=\text{constante}
\quad\Longrightarrow\quad
\vect{a}=0
\]

es falsa en general.

Sólo una \textbf{velocidad vectorial constante} implica:

\[
\vect{a}=0.
\]

Una partícula puede conservar su rapidez y, al mismo tiempo, acelerar porque
cambia la dirección de \(\vect{v}\).
\end{errorfrecuente}


\section{Dirección de la aceleración}

La aceleración no tiene por qué apuntar en la misma dirección que la velocidad.

Podemos encontrar situaciones en las que:

\begin{itemize}
    \item \(\vect{a}\) sea paralela a \(\vect{v}\);
    \item \(\vect{a}\) sea opuesta a \(\vect{v}\);
    \item \(\vect{a}\) forme un ángulo con \(\vect{v}\);
    \item \(\vect{a}\) sea perpendicular a \(\vect{v}\) en cierto instante.
\end{itemize}

\begin{interpretacion}
La velocidad responde principalmente a:

\begin{quote}
¿hacia dónde y qué tan rápido se mueve la partícula?
\end{quote}

La aceleración responde a:

\begin{quote}
¿cómo está cambiando ese vector velocidad?
\end{quote}

Por ello, no existe una regla general que obligue a ambos vectores a tener la
misma dirección.
\end{interpretacion}


\section{Posición, velocidad y aceleración: estructura matemática}

Las relaciones fundamentales pueden resumirse mediante:

\[
\vect{r}(t)
\]

\[
\Downarrow\quad
\deriv{}{t}
\]

\[
\vect{v}(t)
=
\deriv{\vect{r}}{t}
\]

\[
\Downarrow\quad
\deriv{}{t}
\]

\[
\vect{a}(t)
=
\deriv{\vect{v}}{t}.
\]

En componentes:

\[
x(t)
\longrightarrow
v_x(t)
\longrightarrow
a_x(t),
\]

\[
y(t)
\longrightarrow
v_y(t)
\longrightarrow
a_y(t),
\]

\[
z(t)
\longrightarrow
v_z(t)
\longrightarrow
a_z(t).
\]

\begin{interpretacion}
Las derivadas no son únicamente operaciones matemáticas.

Expresan relaciones físicas:

\[
\deriv{x}{t}
\]

mide qué tan rápidamente cambia la coordenada \(x\), mientras:

\[
\deriv{v_x}{t}
\]

mide qué tan rápidamente cambia la componente \(x\) de la velocidad.
\end{interpretacion}


\section{La relación integral}

También podemos recorrer el proceso en sentido contrario.

Si conocemos la aceleración:

\[
\vect{a}(t),
\]

el cambio de velocidad entre \(t_i\) y \(t_f\) es:

\[
\boxed{
\Delta\vect{v}
=
\int_{t_i}^{t_f}
\vect{a}(t)\,\dd t
}.
\]

Y si conocemos la velocidad:

\[
\boxed{
\Delta\vect{r}
=
\int_{t_i}^{t_f}
\vect{v}(t)\,\dd t
}.
\]

Por componentes:

\[
\Delta x
=
\int_{t_i}^{t_f}
v_x(t)\,\dd t,
\]

\[
\Delta y
=
\int_{t_i}^{t_f}
v_y(t)\,\dd t,
\]

\[
\Delta z
=
\int_{t_i}^{t_f}
v_z(t)\,\dd t.
\]

\begin{interpretacion}
La integración representa acumulación.

Acumular pequeños cambios de velocidad produce un cambio total de velocidad.

Acumular pequeños desplazamientos produce el desplazamiento total.
\end{interpretacion}


\section{El método de componentes rectangulares}

El procedimiento más general para trabajar con movimiento en dos o tres
dimensiones consiste en escribir todas las magnitudes vectoriales mediante sus
componentes rectangulares.

Por ejemplo:

\[
\vect{r}
=
x\ihat+y\jhat,
\]

\[
\vect{v}
=
v_x\ihat+v_y\jhat,
\]

\[
\vect{a}
=
a_x\ihat+a_y\jhat.
\]

Las ecuaciones vectoriales pueden entonces convertirse en ecuaciones para cada
eje.

\begin{estrategia}[Idea central]
Para analizar un movimiento multidimensional:

\begin{enumerate}
    \item defina el sistema de referencia;
    \item escriba los vectores mediante componentes;
    \item aplique las relaciones cinemáticas a cada componente;
    \item utilice el mismo tiempo para todas las componentes;
    \item reconstruya el vector final;
    \item determine magnitud y dirección sólo después de obtener las componentes.
\end{enumerate}
\end{estrategia}

Este procedimiento es preferible a intentar trabajar directamente con
magnitudes y ángulos durante todo el cálculo.


\section{¿Son independientes los movimientos en cada eje?}

Con ejes cartesianos fijos, las ecuaciones vectoriales pueden escribirse por
componentes.

Por ejemplo:

\[
\vect{a}
=
a_x\ihat+a_y\jhat
\]

equivale a estudiar:

\[
a_x,
\qquad
a_y.
\]

Esto permite resolver las ecuaciones correspondientes a \(x\) y \(y\) por
separado.

Sin embargo, no significa que estemos estudiando dos partículas distintas.

\begin{interpretacion}
Las componentes comparten una condición fundamental:

\[
\boxed{\text{el mismo tiempo }t}.
\]

La partícula posee simultáneamente una posición \(x(t)\) y una posición \(y(t)\).

La trayectoria completa aparece al combinar ambas componentes.
\end{interpretacion}

Además, algunos sistemas físicos pueden imponer restricciones que relacionen
las diferentes componentes. Por ello, la posibilidad de escribir las ecuaciones
por componentes no debe confundirse con una independencia física universal.


\section{Ejemplo: de la posición a la trayectoria}

\begin{ejemplo}[Partícula con movimiento curvilíneo]
Una partícula se mueve en el plano según:

\[
\vect{r}(t)
=
(1+2t)\ihat
+
(4t-t^2)\jhat,
\]

donde las posiciones se expresan en metros y el tiempo en segundos.

Determine:

\begin{enumerate}[label=\alph*)]
    \item las ecuaciones \(x(t)\) y \(y(t)\);
    \item la velocidad;
    \item la aceleración;
    \item la posición y la velocidad en \(t=\SI{1.0}{\second}\);
    \item la ecuación de la trayectoria.
\end{enumerate}
\end{ejemplo}

Las componentes de posición son:

\[
x(t)=1+2t,
\]

\[
y(t)=4t-t^2.
\]

La velocidad es:

\[
\vect{v}
=
\deriv{\vect{r}}{t}.
\]

Por componentes:

\[
v_x
=
\deriv{}{t}(1+2t)
=
2,
\]

\[
v_y
=
\deriv{}{t}(4t-t^2)
=
4-2t.
\]

Por tanto:

\[
\boxed{
\vect{v}(t)
=
2\ihat+(4-2t)\jhat
}
\]

en \(\si{\meter\per\second}\).

La aceleración es:

\[
\vect{a}
=
\deriv{\vect{v}}{t}.
\]

Entonces:

\[
a_x=0,
\]

\[
a_y=-2.
\]

Así:

\[
\boxed{
\vect{a}
=
-2\jhat\,
\si{\meter\per\second\squared}
}.
\]

Para:

\[
t=\SI{1.0}{\second},
\]

la posición es:

\[
x=1+2(1)=\SI{3}{\meter},
\]

\[
y=4(1)-(1)^2=\SI{3}{\meter}.
\]

Por tanto:

\[
\boxed{
\vect{r}
=
(3\ihat+3\jhat)\,\si{\meter}
}.
\]

La velocidad en ese instante es:

\[
\vect{v}
=
(2\ihat+2\jhat)\,
\si{\meter\per\second}.
\]

La rapidez es:

\[
v
=
\sqrt{2^2+2^2}\,
\si{\meter\per\second}.
\]

Así:

\[
\boxed{
v
=
2\sqrt{2}\,\si{\meter\per\second}
\approx
\SI{2.83}{\meter\per\second}
}.
\]

Como ambas componentes de la velocidad son positivas e iguales:

\[
\theta=45\grado.
\]

Para obtener la trayectoria despejamos el tiempo de:

\[
x=1+2t.
\]

Entonces:

\[
t=\frac{x-1}{2}.
\]

Sustituyendo en:

\[
y=4t-t^2,
\]

obtenemos:

\[
y
=
2(x-1)
-
\frac{(x-1)^2}{4}.
\]

Por tanto:

\[
\boxed{
y
=
2(x-1)
-
\frac{(x-1)^2}{4}
}.
\]

\begin{interpretacion}
La partícula tiene:

\[
v_x=\SI{2}{\meter\per\second}
\]

constante, mientras que \(v_y\) cambia con el tiempo.

Por ello, aunque una componente del movimiento sea uniforme, la trayectoria
completa no tiene por qué ser rectilínea.
\end{interpretacion}


\section{Movimiento con aceleración constante en varias dimensiones}

Consideremos ahora el caso en el que el vector aceleración permanece constante:

\[
\boxed{
\vect{a}=\text{constante}
}.
\]

Esto significa que todas sus componentes cartesianas son constantes:

\[
a_x=\text{constante},
\]

\[
a_y=\text{constante},
\]

\[
a_z=\text{constante}.
\]

\begin{modelo}[Partícula con aceleración constante]
Durante el intervalo considerado suponemos que el vector aceleración mantiene:

\begin{itemize}
    \item la misma magnitud;
    \item la misma dirección;
    \item el mismo sentido.
\end{itemize}

Esta condición permite integrar directamente las ecuaciones cinemáticas.
\end{modelo}


\section{Velocidad para aceleración constante}

Partimos de:

\[
\vect{a}
=
\deriv{\vect{v}}{t}.
\]

Si \(\vect{a}\) es constante:

\[
\dd\vect{v}
=
\vect{a}\,\dd t.
\]

Integramos desde:

\[
t=0
\]

hasta un instante \(t\), y desde:

\[
\vect{v}_0
\]

hasta:

\[
\vect{v}.
\]

Entonces:

\[
\int_{\vect{v}_0}^{\vect{v}}
\dd\vect{v}
=
\int_0^t
\vect{a}\,\dd t.
\]

Por tanto:

\[
\vect{v}-\vect{v}_0
=
\vect{a}t.
\]

Finalmente:

\[
\boxed{
\vect{v}
=
\vect{v}_0+\vect{a}t
}.
\]

\begin{interpretacion}
El término:

\[
\vect{a}t
\]

representa el cambio de velocidad producido durante el intervalo.

La ecuación es vectorial. Por tanto, no afirma que \(\vect{v}\),
\(\vect{v}_0\) y \(\vect{a}\) tengan necesariamente la misma dirección.
\end{interpretacion}


\subsection{Forma por componentes}

En tres dimensiones:

\[
\boxed{
v_x=v_{0x}+a_xt
},
\]

\[
\boxed{
v_y=v_{0y}+a_yt
},
\]

\[
\boxed{
v_z=v_{0z}+a_zt
}.
\]


\section{Posición para aceleración constante}

Sabemos que:

\[
\vect{v}
=
\deriv{\vect{r}}{t}.
\]

Además:

\[
\vect{v}
=
\vect{v}_0+\vect{a}t.
\]

Por tanto:

\[
\deriv{\vect{r}}{t}
=
\vect{v}_0+\vect{a}t.
\]

Multiplicando por \(\dd t\):

\[
\dd\vect{r}
=
(\vect{v}_0+\vect{a}t)\,\dd t.
\]

Integramos desde la posición inicial \(\vect{r}_0\) hasta \(\vect{r}\):

\[
\vect{r}-\vect{r}_0
=
\int_0^t
(\vect{v}_0+\vect{a}t)\,\dd t.
\]

Como \(\vect{v}_0\) y \(\vect{a}\) son constantes:

\[
\vect{r}-\vect{r}_0
=
\vect{v}_0t+\frac{1}{2}\vect{a}t^2.
\]

Así:

\[
\boxed{
\vect{r}
=
\vect{r}_0
+
\vect{v}_0t
+
\frac{1}{2}\vect{a}t^2
}.
\]

Equivalentemente:

\[
\boxed{
\Delta\vect{r}
=
\vect{v}_0t
+
\frac{1}{2}\vect{a}t^2
}.
\]


\subsection{Interpretación de los términos}

La ecuación:

\[
\vect{r}
=
\vect{r}_0
+
\vect{v}_0t
+
\frac{1}{2}\vect{a}t^2
\]

puede leerse físicamente.

El término:

\[
\vect{r}_0
\]

representa la posición inicial.

El término:

\[
\vect{v}_0t
\]

representa el desplazamiento que tendría la partícula si mantuviera su velocidad
inicial.

El término:

\[
\frac{1}{2}\vect{a}t^2
\]

representa la modificación del desplazamiento producida por la aceleración.


\section{Ecuaciones por componentes para aceleración constante}

La ecuación vectorial de posición equivale a tres ecuaciones escalares.

Para \(x\):

\[
\boxed{
x
=
x_0+v_{0x}t+\frac{1}{2}a_xt^2
}.
\]

Para \(y\):

\[
\boxed{
y
=
y_0+v_{0y}t+\frac{1}{2}a_yt^2
}.
\]

Para \(z\):

\[
\boxed{
z
=
z_0+v_{0z}t+\frac{1}{2}a_zt^2
}.
\]

De forma similar:

\[
v_x=v_{0x}+a_xt,
\]

\[
v_y=v_{0y}+a_yt,
\]

\[
v_z=v_{0z}+a_zt.
\]

\begin{idea}[Una sola ecuación vectorial, varias ecuaciones escalares]
Las ecuaciones por componentes no son nuevas leyes.

Son la expresión cartesiana de las mismas ecuaciones vectoriales:

\[
\vect{v}
=
\vect{v}_0+\vect{a}t,
\]

\[
\vect{r}
=
\vect{r}_0+\vect{v}_0t+\frac{1}{2}\vect{a}t^2.
\]
\end{idea}


\section{Caso especial: una componente de aceleración nula}

Supongamos que:

\[
a_x=0.
\]

Entonces:

\[
v_x=v_{0x},
\]

es decir, la componente horizontal de la velocidad permanece constante.

Además:

\[
x=x_0+v_{0x}t.
\]

Esto no implica que toda la velocidad sea constante.

Si:

\[
a_y\neq0,
\]

la componente \(v_y\) continúa cambiando.

\begin{interpretacion}
Una partícula puede tener simultáneamente:

\[
\text{movimiento uniforme en una coordenada}
\]

y:

\[
\text{movimiento acelerado en otra}.
\]

El movimiento completo se obtiene combinando ambas componentes en cada instante.
\end{interpretacion}


\section{Caso especial: aceleración nula}

Si:

\[
\vect{a}=\vect{0},
\]

entonces:

\[
\vect{v}=\vect{v}_0=\text{constante}.
\]

La posición queda:

\[
\boxed{
\vect{r}
=
\vect{r}_0+\vect{v}_0t
}.
\]

Como la dirección de \(\vect{v}_0\) permanece constante, la trayectoria es una
línea recta.

\begin{advertencia}
No debe confundirse:

\[
\vect{a}=\vect{0}
\]

con:

\[
v=\text{constante}.
\]

Aceleración nula exige que el \textbf{vector velocidad completo} sea constante,
no solamente su magnitud.
\end{advertencia}


\section{Una relación sin el tiempo}

En una dimensión obtuvimos una relación entre velocidad, aceleración y
desplazamiento que no contiene explícitamente el tiempo.

Existe una generalización útil para aceleración vectorial constante.

De:

\[
\vect{v}
=
\vect{v}_0+\vect{a}t
\]

obtenemos:

\[
v^2
=
(\vect{v}_0+\vect{a}t)
\pdot
(\vect{v}_0+\vect{a}t).
\]

Entonces:

\[
v^2
=
v_0^2
+
2\vect{a}\pdot\vect{v}_0t
+
a^2t^2.
\]

Por otra parte:

\[
\Delta\vect{r}
=
\vect{v}_0t
+
\frac{1}{2}\vect{a}t^2.
\]

Multiplicando escalarmente por \(2\vect{a}\):

\[
2\vect{a}\pdot\Delta\vect{r}
=
2\vect{a}\pdot\vect{v}_0t
+
a^2t^2.
\]

Comparando ambas expresiones:

\[
\boxed{
v^2
=
v_0^2
+
2\vect{a}\pdot\Delta\vect{r}
}.
\]

\begin{condicion}
Esta relación requiere:

\[
\vect{a}=\text{constante}.
\]

Además, \(v\) y \(v_0\) representan las magnitudes de las velocidades.

La ecuación permite relacionar rapidez y desplazamiento, pero por sí sola no
determina toda la dirección del vector velocidad final.
\end{condicion}


\section{Ejemplo: aceleración constante en un plano}

\begin{ejemplo}[Movimiento bidimensional]
Una partícula tiene inicialmente:

\[
\vect{r}_0
=
(-2\ihat+\jhat)\,\si{\meter},
\]

\[
\vect{v}_0
=
(3\ihat+2\jhat)\,
\si{\meter\per\second},
\]

y experimenta la aceleración constante:

\[
\vect{a}
=
(\ihat-2\jhat)\,
\si{\meter\per\second\squared}.
\]

Determine su posición, velocidad y rapidez en:

\[
t=\SI{3.0}{\second}.
\]
\end{ejemplo}

Primero analizamos las componentes.

Los datos son:

\[
x_0=-\SI{2}{\meter},
\qquad
y_0=\SI{1}{\meter},
\]

\[
v_{0x}=\SI{3}{\meter\per\second},
\qquad
v_{0y}=\SI{2}{\meter\per\second},
\]

\[
a_x=\SI{1}{\meter\per\second\squared},
\qquad
a_y=-\SI{2}{\meter\per\second\squared}.
\]

Para la componente \(x\):

\[
x
=
x_0+v_{0x}t+\frac{1}{2}a_xt^2.
\]

Entonces:

\[
x
=
-2+3(3)+\frac{1}{2}(1)(3)^2.
\]

Así:

\[
x=\SI{11.5}{\meter}.
\]

Para la componente \(y\):

\[
y
=
y_0+v_{0y}t+\frac{1}{2}a_yt^2.
\]

Entonces:

\[
y
=
1+2(3)+\frac{1}{2}(-2)(3)^2.
\]

Por tanto:

\[
y=-\SI{2.0}{\meter}.
\]

La posición es:

\[
\boxed{
\vect{r}
=
(11.5\ihat-2.0\jhat)\,\si{\meter}
}.
\]

Ahora calculamos la velocidad.

Para \(x\):

\[
v_x
=
v_{0x}+a_xt.
\]

Entonces:

\[
v_x
=
3+1(3)
=
\SI{6.0}{\meter\per\second}.
\]

Para \(y\):

\[
v_y
=
v_{0y}+a_yt.
\]

Entonces:

\[
v_y
=
2-2(3)
=
-\SI{4.0}{\meter\per\second}.
\]

Así:

\[
\boxed{
\vect{v}
=
(6.0\ihat-4.0\jhat)\,
\si{\meter\per\second}
}.
\]

La rapidez es:

\[
v
=
\sqrt{6.0^2+(-4.0)^2}.
\]

Por tanto:

\[
\boxed{
v
=
\sqrt{52}\,
\si{\meter\per\second}
\approx
\SI{7.21}{\meter\per\second}
}.
\]

Como:

\[
v_x>0
\]

y:

\[
v_y<0,
\]

la velocidad apunta hacia el cuarto cuadrante.

El ángulo respecto del eje \(+x\) tiene magnitud:

\[
\theta
=
\arctan\left(\frac{4}{6}\right)
\approx
33.7\grado.
\]

Por tanto, la velocidad apunta aproximadamente:

\[
\boxed{
33.7\grado
\text{ por debajo del eje }+x
}.
\]

\begin{interpretacion}
La aceleración no es paralela a la velocidad.

Por ello, el vector velocidad no sólo puede modificar su magnitud, sino también
su dirección durante el movimiento.
\end{interpretacion}


\section{Movimiento en tres dimensiones}

No aparece un principio cinemático nuevo al pasar de dos a tres dimensiones.

Simplemente añadimos una tercera componente.

La posición es:

\[
\vect{r}
=
x\ihat+y\jhat+z\khat.
\]

La velocidad:

\[
\vect{v}
=
v_x\ihat+v_y\jhat+v_z\khat.
\]

La aceleración:

\[
\vect{a}
=
a_x\ihat+a_y\jhat+a_z\khat.
\]

Sus magnitudes son:

\[
r
=
\sqrt{x^2+y^2+z^2},
\]

\[
v
=
\sqrt{v_x^2+v_y^2+v_z^2},
\]

\[
a
=
\sqrt{a_x^2+a_y^2+a_z^2}.
\]

\begin{interpretacion}
La dificultad conceptual del movimiento tridimensional no está en una nueva
definición de velocidad o aceleración.

Está principalmente en organizar correctamente las tres componentes y
reconstruir después el vector físico completo.
\end{interpretacion}


\section{Ejemplo: movimiento espacial}

\begin{ejemplo}[Movimiento de un dron]
La posición de un dron se modela mediante:

\[
\vect{r}(t)
=
(1+t)\ihat
+
2t^2\jhat
+
\left(3-\frac{1}{2}t\right)\khat,
\]

con posiciones en metros y tiempo en segundos.

Determine la posición, velocidad, aceleración y rapidez en:

\[
t=\SI{2.0}{\second}.
\]
\end{ejemplo}

Las componentes de posición son:

\[
x=1+t,
\]

\[
y=2t^2,
\]

\[
z=3-\frac{1}{2}t.
\]

En:

\[
t=\SI{2}{\second},
\]

obtenemos:

\[
x=\SI{3}{\meter},
\]

\[
y=\SI{8}{\meter},
\]

\[
z=\SI{2}{\meter}.
\]

Por tanto:

\[
\boxed{
\vect{r}(2)
=
(3\ihat+8\jhat+2\khat)\,\si{\meter}
}.
\]

La velocidad es:

\[
\vect{v}
=
\deriv{\vect{r}}{t}.
\]

Entonces:

\[
\vect{v}(t)
=
\ihat+4t\jhat-\frac{1}{2}\khat
\]

en \(\si{\meter\per\second}\).

Para \(t=\SI{2}{\second}\):

\[
\boxed{
\vect{v}(2)
=
\left(
\ihat+8\jhat-\frac{1}{2}\khat
\right)
\si{\meter\per\second}
}.
\]

La aceleración es:

\[
\vect{a}
=
\deriv{\vect{v}}{t}.
\]

Por tanto:

\[
\boxed{
\vect{a}
=
4\jhat\,
\si{\meter\per\second\squared}
}.
\]

La rapidez en \(t=\SI{2}{\second}\) es:

\[
v
=
\sqrt{
1^2
+
8^2
+
\left(-\frac{1}{2}\right)^2
}.
\]

Así:

\[
v
=
\sqrt{65.25}\,
\si{\meter\per\second}.
\]

Por tanto:

\[
\boxed{
v\approx
\SI{8.08}{\meter\per\second}
}.
\]

\begin{interpretacion}
Aunque el movimiento ocurre en tres dimensiones, el procedimiento no cambió:

\begin{enumerate}
    \item identificamos las componentes;
    \item derivamos cada componente;
    \item reconstruimos los vectores;
    \item calculamos la magnitud sólo al final.
\end{enumerate}
\end{interpretacion}


\section{Representaciones gráficas del movimiento multidimensional}

En una dimensión una gráfica \(x-t\) permite representar directamente la
posición en función del tiempo.

En varias dimensiones conviene distinguir diferentes representaciones.


\subsection{Gráfica \(x-t\)}

Representa:

\[
x=x(t).
\]

Su pendiente instantánea es:

\[
v_x=\deriv{x}{t}.
\]


\subsection{Gráfica \(y-t\)}

Representa:

\[
y=y(t).
\]

Su pendiente es:

\[
v_y=\deriv{y}{t}.
\]


\subsection{Gráfica \(z-t\)}

Representa:

\[
z=z(t).
\]

Su pendiente es:

\[
v_z=\deriv{z}{t}.
\]


\subsection{Gráfica de la trayectoria}

Una gráfica con:

\[
x
\]

en un eje espacial y:

\[
y
\]

en el otro representa la trayectoria en el plano.

En esta gráfica el tiempo no aparece necesariamente como uno de los ejes.

\begin{errorfrecuente}[Trayectoria y gráfica posición--tiempo]
Una gráfica \(x-t\) no debe confundirse con una trayectoria.

En una gráfica \(x-t\):

\begin{itemize}
    \item un eje representa posición;
    \item el otro representa tiempo.
\end{itemize}

En una trayectoria \(y-x\):

\begin{itemize}
    \item ambos ejes representan coordenadas espaciales.
\end{itemize}

Son representaciones físicamente diferentes.
\end{errorfrecuente}


\section{Cómo reconstruir el movimiento a partir de sus componentes}

Supongamos que conocemos:

\[
x(t)
\]

y:

\[
y(t).
\]

Para comprender el movimiento completo podemos seguir esta secuencia:

\begin{enumerate}
    \item calcular la posición:
    \[
    \vect{r}(t)
    =
    x(t)\ihat+y(t)\jhat;
    \]

    \item derivar para obtener:
    \[
    v_x(t),
    \qquad
    v_y(t);
    \]

    \item formar:
    \[
    \vect{v}(t)
    =
    v_x(t)\ihat+v_y(t)\jhat;
    \]

    \item calcular:
    \[
    v(t)
    =
    \sqrt{v_x^2+v_y^2};
    \]

    \item derivar nuevamente:
    \[
    a_x(t),
    \qquad
    a_y(t);
    \]

    \item formar:
    \[
    \vect{a}(t)
    =
    a_x(t)\ihat+a_y(t)\jhat;
    \]

    \item eliminar \(t\), cuando sea útil, para obtener la trayectoria:
    \[
    y=y(x).
    \]
\end{enumerate}

Esta secuencia conecta las representaciones:

\[
\text{paramétrica}
\longleftrightarrow
\text{vectorial}
\longleftrightarrow
\text{geométrica}.
\]


\section{Estrategia general para resolver problemas}

\begin{estrategia}[Movimiento en dos y tres dimensiones]
Ante un problema de cinemática multidimensional:

\begin{enumerate}
    \item Identifique el objeto que será modelado como partícula.

    \item Determine si el movimiento requiere una, dos o tres dimensiones.

    \item Defina un sistema de referencia:
    \[
    O,\qquad x,\qquad y,\qquad z.
    \]

    \item Establezca claramente los sentidos positivos.

    \item Dibuje la trayectoria o un esquema cuando facilite la interpretación.

    \item Exprese las posiciones, velocidades y aceleraciones mediante componentes
    rectangulares.

    \item Determine qué modelo cinemático corresponde:
    \begin{itemize}
        \item velocidad constante;
        \item aceleración constante;
        \item funciones generales de \(t\).
    \end{itemize}

    \item Escriba las relaciones físicas por componentes.

    \item Utilice el mismo tiempo en todos los ejes.

    \item Resuelva primero las componentes desconocidas.

    \item Reconstruya después el vector completo.

    \item Calcule magnitud y dirección cuando sean necesarias.

    \item Compruebe:
    \begin{itemize}
        \item unidades;
        \item signos;
        \item cuadrante;
        \item sentido físico;
        \item coherencia con la trayectoria.
    \end{itemize}
\end{enumerate}
\end{estrategia}


\section{Selección del método}

No todos los problemas multidimensionales requieren el mismo procedimiento.

\subsection{Si se conoce \texorpdfstring{\(\vect{r}(t)\)}{r(t)}}

Utilice:

\[
\vect{v}
=
\deriv{\vect{r}}{t}
\]

y:

\[
\vect{a}
=
\deriv{\vect{v}}{t}.
\]


\subsection{Si se conoce \texorpdfstring{\(\vect{a}(t)\)}{a(t)}}

Puede obtener el cambio de velocidad mediante:

\[
\Delta\vect{v}
=
\int\vect{a}\,\dd t.
\]

Después, a partir de \(\vect{v}(t)\), puede determinarse la posición.


\subsection{Si la aceleración es constante}

Utilice directamente:

\[
\vect{v}
=
\vect{v}_0+\vect{a}t
\]

y:

\[
\vect{r}
=
\vect{r}_0
+
\vect{v}_0t
+
\frac{1}{2}\vect{a}t^2.
\]

Luego trabaje por componentes.


\section{Errores frecuentes}

\begin{errorfrecuente}
\textbf{Sumar magnitudes en lugar de componentes.}

Si:

\[
\vect{A}=A_x\ihat+A_y\jhat
\]

y:

\[
\vect{B}=B_x\ihat+B_y\jhat,
\]

la suma correcta es:

\[
\vect{A}+\vect{B}
=
(A_x+B_x)\ihat
+
(A_y+B_y)\jhat.
\]

No deben sumarse directamente las magnitudes salvo en casos geométricos
particulares.
\end{errorfrecuente}


\begin{errorfrecuente}
\textbf{Confundir el vector de posición con el desplazamiento.}

El vector:

\[
\vect{r}
\]

describe dónde está la partícula respecto del origen.

El vector:

\[
\Delta\vect{r}
\]

describe el cambio entre dos posiciones.
\end{errorfrecuente}


\begin{errorfrecuente}
\textbf{Confundir desplazamiento con distancia recorrida.}

El desplazamiento depende sólo de los extremos.

La distancia depende de toda la trayectoria.
\end{errorfrecuente}


\begin{errorfrecuente}
\textbf{Creer que la velocidad siempre apunta en la dirección del vector de
posición.}

La velocidad instantánea es tangente a la trayectoria.

No tiene por qué apuntar radialmente desde el origen.
\end{errorfrecuente}


\begin{errorfrecuente}
\textbf{Creer que la aceleración siempre apunta en la dirección del movimiento.}

La aceleración indica cómo cambia \(\vect{v}\).

Puede tener una dirección diferente de la velocidad.
\end{errorfrecuente}


\begin{errorfrecuente}
\textbf{Concluir que rapidez constante significa aceleración cero.}

Una rapidez constante sólo implica que:

\[
\norm{\vect{v}}
\]

no cambia.

Si cambia la dirección de \(\vect{v}\), existe aceleración.
\end{errorfrecuente}


\begin{errorfrecuente}
\textbf{Usar tiempos distintos para cada componente.}

Las ecuaciones:

\[
x(t),
\qquad
y(t),
\qquad
z(t)
\]

describen simultáneamente la misma partícula.

El tiempo es común a todas las componentes.
\end{errorfrecuente}


\begin{errorfrecuente}
\textbf{Aplicar las ecuaciones de aceleración constante cuando
\(\vect{a}\) no es constante.}

Las relaciones:

\[
\vect{v}
=
\vect{v}_0+\vect{a}t
\]

y:

\[
\vect{r}
=
\vect{r}_0+\vect{v}_0t+\frac{1}{2}\vect{a}t^2
\]

requieren aceleración vectorial constante durante el intervalo considerado.
\end{errorfrecuente}


\begin{errorfrecuente}
\textbf{Obtener un ángulo con la calculadora y no comprobar el cuadrante.}

La relación:

\[
\tan\theta=\frac{A_y}{A_x}
\]

no determina por sí sola cuál de los posibles cuadrantes contiene el vector.

Siempre deben analizarse los signos de las componentes.
\end{errorfrecuente}


\section{Aplicaciones}

La descripción vectorial del movimiento aparece en numerosos sistemas reales.

\begin{aplicacion}[Navegación y seguimiento]
Los sistemas de navegación pueden representar la posición de un vehículo
mediante coordenadas que cambian con el tiempo.

A partir de esas posiciones pueden determinarse desplazamientos y velocidades.

Cuando el movimiento ocurre sobre una superficie se utilizan esencialmente dos
coordenadas; para vehículos aéreos se requiere también información sobre la
altura.
\end{aplicacion}


\begin{aplicacion}[Robótica]
Un robot móvil puede programarse mediante coordenadas de posición:

\[
(x,y)
\]

o:

\[
(x,y,z).
\]

La planificación de su movimiento exige controlar cómo cambian esas coordenadas
con el tiempo y, por tanto, cómo cambian sus vectores de velocidad y
aceleración.
\end{aplicacion}


\begin{aplicacion}[Registro experimental de trayectorias]
Mediante una cámara puede registrarse la posición de un objeto en diferentes
instantes.

Si se establece una escala espacial y un sistema de coordenadas, pueden
obtenerse datos:

\[
x(t),
\qquad
y(t).
\]

A partir de ellos pueden estimarse las componentes de velocidad y aceleración.

Un intervalo temporal cada vez menor permite aproximarse mejor a las magnitudes
instantáneas, aunque en una medición real aparecen limitaciones experimentales.
\end{aplicacion}


\section{Secuencia sugerida para impartir la clase}

\begin{notapedagogica}
Una secuencia posible para desarrollar este capítulo directamente en clase es:

\begin{enumerate}
    \item Dibujar primero una carretera recta y preguntar cuántas coordenadas se
    necesitan para localizar un automóvil.

    \item Dibujar después una trayectoria curva sobre un plano y preguntar por qué
    una sola coordenada deja de ser suficiente.

    \item Recordar brevemente el método de componentes rectangulares.

    \item Introducir el vector de posición:
    \[
    \vect{r}=x\ihat+y\jhat.
    \]

    \item Mostrar que diferentes valores de \(t\) producen diferentes posiciones y
    construyen una trayectoria.

    \item Diferenciar cuidadosamente:
    \[
    \vect{r},
    \qquad
    \Delta\vect{r},
    \qquad
    d.
    \]

    \item Construir la velocidad media a partir del desplazamiento.

    \item Reducir conceptualmente el intervalo temporal hasta introducir:
    \[
    \vect{v}=\deriv{\vect{r}}{t}.
    \]

    \item Utilizar un dibujo para mostrar que la velocidad instantánea es tangente
    a la trayectoria.

    \item Obtener:
    \[
    v_x=\deriv{x}{t},
    \qquad
    v_y=\deriv{y}{t}.
    \]

    \item Diferenciar velocidad y rapidez.

    \item Introducir el cambio del vector velocidad y construir:
    \[
    \vect{a}=\deriv{\vect{v}}{t}.
    \]

    \item Mostrar mediante dos vectores de igual magnitud y distinta dirección que
    puede existir aceleración con rapidez constante.

    \item Resolver un ejemplo a partir de una función:
    \[
    \vect{r}(t).
    \]

    \item Introducir después el modelo de aceleración constante.

    \item Deducir primero:
    \[
    \vect{v}=\vect{v}_0+\vect{a}t,
    \]
    y después:
    \[
    \vect{r}
    =
    \vect{r}_0+\vect{v}_0t+\frac{1}{2}\vect{a}t^2.
    \]

    \item Descomponer explícitamente ambas ecuaciones en componentes \(x\), \(y\)
    y \(z\).

    \item Resolver un ejemplo completo mediante componentes rectangulares.

    \item Finalizar comparando:
    \[
    x-t,\qquad y-t,\qquad y-x,
    \]
    para evitar que los estudiantes confundan una gráfica temporal con una
    trayectoria.
\end{enumerate}
\end{notapedagogica}


\section{Preguntas de comprensión}

\begin{enumerate}
    \item ¿Qué hace que un movimiento sea bidimensional?

    \item ¿Un cuerpo tridimensional necesariamente realiza un movimiento
    tridimensional?

    \item ¿Qué representa físicamente el vector de posición?

    \item ¿Qué diferencia existe entre:
    \[
    \vect{r}
    \]
    y:
    \[
    \Delta\vect{r}?
    \]

    \item ¿La magnitud del desplazamiento siempre coincide con la distancia
    recorrida?

    \item Si una partícula regresa a su punto inicial después de seguir una
    trayectoria curva, ¿cuál es su desplazamiento?

    \item ¿Qué dirección posee la velocidad media respecto del desplazamiento?

    \item ¿Por qué la velocidad instantánea es tangente a la trayectoria?

    \item ¿Puede el vector velocidad ser perpendicular al vector de posición?

    \item ¿Qué representan físicamente \(v_x\), \(v_y\) y \(v_z\)?

    \item ¿Qué diferencia existe entre velocidad y rapidez?

    \item ¿Puede una partícula tener rapidez constante y aceleración distinta de
    cero?

    \item ¿Qué debe cambiar necesariamente para que exista aceleración?

    \item ¿La aceleración debe tener siempre la misma dirección que la velocidad?

    \item Si:
    \[
    a_x=0,
    \]
    ¿qué puede afirmarse acerca de \(v_x\)?

    \item Si:
    \[
    \vect{a}=\vect{0},
    \]
    ¿qué puede afirmarse acerca del vector velocidad?

    \item ¿Por qué los movimientos descritos por \(x(t)\) y \(y(t)\) deben utilizar
    el mismo tiempo?

    \item ¿Qué diferencia existe entre una gráfica \(x-t\) y una gráfica \(y-x\)?

    \item ¿Qué condición debe cumplirse para utilizar:
    \[
    \vect{v}=\vect{v}_0+\vect{a}t?
    \]

    \item ¿Qué información adicional aporta escribir una respuesta como vector en
    lugar de proporcionar únicamente su magnitud?

    \item Si dos velocidades poseen la misma magnitud pero diferentes direcciones,
    ¿su cambio de velocidad es necesariamente cero?

    \item ¿Por qué el método de componentes rectangulares resulta útil para
    movimientos en varias dimensiones?

    \item ¿Puede una componente de la velocidad permanecer constante mientras otra
    cambia?

    \item ¿Qué información se pierde si solamente conocemos la rapidez de una
    partícula?
\end{enumerate}


\begin{resumen}
El movimiento en dos y tres dimensiones es una extensión vectorial de la
cinemática desarrollada para el movimiento rectilíneo.

La posición de una partícula se describe mediante:

\[
\vect{r}
=
x\ihat+y\jhat+z\khat.
\]

El desplazamiento es el cambio del vector de posición:

\[
\Delta\vect{r}
=
\vect{r}_f-\vect{r}_i.
\]

Su magnitud no debe confundirse con la distancia total recorrida.

La velocidad media es:

\[
\vect{\bar{v}}
=
\frac{\Delta\vect{r}}{\Delta t}.
\]

La velocidad instantánea es la tasa de cambio de la posición:

\[
\vect{v}
=
\deriv{\vect{r}}{t},
\]

y es tangente a la trayectoria.

Sus componentes cartesianas son:

\[
v_x=\deriv{x}{t},
\qquad
v_y=\deriv{y}{t},
\qquad
v_z=\deriv{z}{t}.
\]

La rapidez es la magnitud del vector velocidad:

\[
v=\norm{\vect{v}}.
\]

La aceleración es la tasa de cambio del vector velocidad:

\[
\vect{a}
=
\deriv{\vect{v}}{t}
=
\derivdos{\vect{r}}{t}.
\]

Por ello, existe aceleración si cambia la magnitud de la velocidad, su dirección
o ambas.

Una rapidez constante no implica necesariamente aceleración nula.

Cuando el vector aceleración es constante:

\[
\vect{v}
=
\vect{v}_0+\vect{a}t,
\]

y:

\[
\vect{r}
=
\vect{r}_0
+
\vect{v}_0t
+
\frac{1}{2}\vect{a}t^2.
\]

Estas relaciones vectoriales pueden aplicarse componente a componente:

\[
x,\qquad y,\qquad z.
\]

El procedimiento general consiste en establecer primero el sistema de
referencia, trabajar mediante componentes rectangulares y reconstruir al final
la magnitud, dirección y sentido del vector buscado.

La idea central puede expresarse mediante la cadena:

\[
\boxed{
\vect{r}(t)
\longrightarrow
\vect{v}(t)
=
\deriv{\vect{r}}{t}
\longrightarrow
\vect{a}(t)
=
\deriv{\vect{v}}{t}
}.
\]

La descripción completa de un movimiento multidimensional surge de combinar
correctamente sus componentes espaciales en un mismo instante de tiempo.
\end{resumen}